% ============================================================
% Paper 88: Odd Perfect Numbers: Structure, Constraints,
%           and the Open Problem
% Author: Michael Fuhrmann
% Specialist Mathematics Project, Build 168
% Date: 2026-03-12
% ============================================================
\documentclass[12pt,a4paper]{amsart}

\usepackage{amsmath,amssymb,amsthm,mathtools,enumitem,booktabs,array,hyperref}
\usepackage[utf8]{inputenc}
\usepackage[T1]{fontenc}
\usepackage{geometry}
\geometry{margin=2.5cm}

% ---- Theorem environments ----
\newtheorem{theorem}{Theorem}[section]
\newtheorem{lemma}[theorem]{Lemma}
\newtheorem{corollary}[theorem]{Corollary}
\newtheorem{proposition}[theorem]{Proposition}
\theoremstyle{definition}
\newtheorem{definition}[theorem]{Definition}
\newtheorem{example}[theorem]{Example}
\theoremstyle{remark}
\newtheorem{remark}[theorem]{Remark}
\newtheorem{conjecture}[theorem]{Conjecture}

% ---- Custom operators ----
\DeclareMathOperator{\ord}{ord}
\newcommand{\N}{\mathbb{N}}
\newcommand{\Z}{\mathbb{Z}}
\newcommand{\Q}{\mathbb{Q}}
\newcommand{\R}{\mathbb{R}}

% ============================================================
\begin{document}

\title{Odd Perfect Numbers:\\
Structure, Constraints, and the Open Problem}

\author{Michael Fuhrmann}
\address{Specialist Mathematics Project, Build 168}
\email{specialist-maths@project.local}
\date{2026-03-12}

\begin{abstract}
A positive integer $n$ is called \emph{perfect} if $\sigma(n) = 2n$,
where $\sigma$ denotes the sum-of-divisors function.
While all known perfect numbers are even and arise from Mersenne primes via
the Euclid--Euler theorem, the question of whether any \emph{odd} perfect
number exists has remained open for over two thousand years, making it
arguably the oldest unsolved problem in mathematics.
In this paper we survey the principal structural constraints on a hypothetical
odd perfect number: Euler's factorisation theorem, which forces $n = p^a m^2$
with $p \equiv a \equiv 1 \pmod{4}$, as well as modern lower bounds
($n > 10^{1500}$, Ochem--Rao 2012), lower bounds on the number of distinct
prime factors ($\omega(n) \geq 9$), bounds on the largest and smallest prime
factors, and a variety of divisibility and congruence restrictions.
We also discuss closely related concepts --- quasiperfect, multiply perfect,
and hyperperfect numbers --- and formulate the central conjecture.
\end{abstract}

\maketitle

\tableofcontents

% ============================================================
\section{Introduction}
% ============================================================

The \emph{sum-of-divisors} function $\sigma : \N \to \N$ assigns to each
positive integer $n$ the sum of all its positive divisors:
\[
  \sigma(n) = \sum_{d \mid n} d.
\]
A number is called \emph{perfect} when $\sigma(n) = 2n$, i.e.\ when the
proper divisors of $n$ sum to $n$ itself.
The first few perfect numbers are $6 = 1+2+3$,
$28 = 1+2+4+7+14$,
$496$, and $8128$.

The ancient Greeks, especially Euclid and Nicomachus, were fascinated by
perfect numbers.
Euclid proved that $2^{p-1}(2^p-1)$ is perfect whenever $2^p - 1$ is prime
(a \emph{Mersenne prime}).
Euler later showed that \emph{every} even perfect number has this form.
Thus even perfect numbers are completely classified modulo the open question
of how many Mersenne primes exist.

The situation for \emph{odd} perfect numbers is entirely different.
Despite centuries of effort and modern computational searches reaching
$n > 10^{1500}$, no odd perfect number has ever been found, nor has anyone
proved one cannot exist.

\begin{conjecture}[No Odd Perfect Numbers]
\label{conj:nopn}
There are no odd perfect numbers.
\end{conjecture}

Conjecture~\ref{conj:nopn} is widely believed to be true.
This paper presents the strongest known structural constraints that any
hypothetical odd perfect number must satisfy.

% ============================================================
\section{Arithmetic Background}
% ============================================================

\subsection{The sum-of-divisors function}

\begin{definition}
For $n \in \N$, set
\[
  \sigma(n) = \sum_{d \mid n} d, \qquad
  s(n) = \sigma(n) - n \quad (\text{sum of \emph{proper} divisors}).
\]
\end{definition}

\begin{proposition}[Multiplicativity]
$\sigma$ is multiplicative on coprime integers and satisfies on prime powers:
\[
  \sigma(p^k) = 1 + p + \cdots + p^k = \frac{p^{k+1}-1}{p-1},
\]
and for $\gcd(a,b)=1$: $\sigma(ab) = \sigma(a)\sigma(b)$.
\end{proposition}

\begin{proof}
This follows from the fact that every divisor of $p_1^{k_1}\cdots p_r^{k_r}$
can be written uniquely as $\prod p_i^{e_i}$ with $0 \leq e_i \leq k_i$,
and the sum factors accordingly.
\end{proof}

\begin{definition}
The \emph{abundancy index} of $n$ is $\sigma(n)/n$.
A perfect number satisfies $\sigma(n)/n = 2$.
\end{definition}

\subsection{Classification of natural numbers by abundancy}

\begin{definition}
A positive integer $n$ is:
\begin{itemize}[nosep]
  \item \emph{deficient} if $\sigma(n) < 2n$;
  \item \emph{perfect} if $\sigma(n) = 2n$;
  \item \emph{abundant} if $\sigma(n) > 2n$.
\end{itemize}
\end{definition}

Every prime $p$ is deficient: $\sigma(p)/p = 1 + 1/p < 2$.
Powers of 2 satisfy $\sigma(2^k)/2^k = 2 - 2^{1-k} < 2$, so also deficient.

% ============================================================
\section{Euler's Structure Theorem}
% ============================================================

The foundational result in the theory of odd perfect numbers is due to Euler:

\begin{theorem}[Euler's Structure Theorem, 1747]
\label{thm:euler}
If $n$ is an odd perfect number, then
\[
  n = p^a \cdot m^2,
\]
where $p$ is a prime, $\gcd(p, m) = 1$, $p \equiv 1 \pmod{4}$,
and $a \equiv 1 \pmod{4}$.
\end{theorem}

\begin{proof}[Proof sketch]
Write $n = p_1^{a_1} \cdots p_r^{a_r}$ with all $p_i$ odd.
Since $\sigma(n) = 2n$, we have
\[
  \prod_{i=1}^r \sigma(p_i^{a_i}) = 2 \prod_{i=1}^r p_i^{a_i}.
\]
The left side must be divisible by exactly one factor of $2$ (since $n$ is
odd and $2n$ contributes a single factor of $2$ beyond $n$).
Now $\sigma(p^a) = 1 + p + \cdots + p^a$.
If $p$ is odd, then $\sigma(p^a)$ is odd if and only if $a$ is even
(since each term is odd and there are $a+1$ of them, so $\sigma(p^a) \equiv a+1 \pmod{2}$).
Thus exactly one exponent $a_j$ must be odd; call this prime $p = p_j$ and its
exponent $a = a_j$.
All other $a_i$ are even, making $n = p^a \cdot m^2$.

Next, $\sigma(p^a) = \frac{p^{a+1}-1}{p-1}$ must be odd (otherwise two factors
contribute the needed factor of $2$).
Checking $\pmod{4}$: since $p$ is odd, either $p \equiv 1$ or $p \equiv 3 \pmod{4}$.
If $p \equiv 3 \pmod{4}$, then $p^{a+1} \equiv (-1)^{a+1} \pmod{4}$.
For $\sigma(p^a)$ to be odd one checks that we need $p \equiv 1 \pmod{4}$
and $a \equiv 1 \pmod{4}$.
\end{proof}

\begin{remark}
The prime $p$ appearing in Euler's theorem is called the \emph{Euler prime}
or \emph{special prime} of $n$.
It satisfies $p \equiv 1 \pmod{4}$ and appears to an odd power $a \equiv 1 \pmod 4$,
so the smallest possibility is $a = 1$ and $p = 5$.
\end{remark}

\begin{corollary}
\label{cor:form}
Any odd perfect number has the form $n \equiv 1 \pmod{4}$.
\end{corollary}

\begin{proof}
$n = p^a m^2$ with $p \equiv 1 \pmod{4}$ and $m$ odd,
so $m^2 \equiv 1 \pmod{4}$ and $p^a \equiv 1 \pmod{4}$, hence $n \equiv 1 \pmod{4}$.
\end{proof}

% ============================================================
\section{Bounds on the Size of an Odd Perfect Number}
% ============================================================

\begin{theorem}[Ochem--Rao, 2012]
\label{thm:ochem}
If $n$ is an odd perfect number, then $n > 10^{1500}$.
\end{theorem}

This remarkable bound was established by Ochem and Rao using a combination
of algebraic constraints and computer-assisted case analysis.
Prior bounds had progressed from Euler's work through successive improvements:
\begin{itemize}[nosep]
  \item Kanold (1941): $n > 10^{20}$
  \item Brent--Cohen--te Riele (1991): $n > 10^{300}$
  \item Ochem--Rao (2012): $n > 10^{1500}$
\end{itemize}

\begin{theorem}[Upper bound on specific prime]
The Euler prime $p$ of an odd perfect number satisfies $p < n^{1/3}$,
and more precisely $p^a < n^{1/2}$ unless $n$ has a very specific form.
\end{theorem}

\begin{remark}
The size bound $n > 10^{1500}$ implies that any odd perfect number would have
at least $1500$ decimal digits, placing it far beyond any conceivable
computational verification.
\end{remark}

% ============================================================
\section{Bounds on the Number of Prime Factors}
% ============================================================

\begin{theorem}[Chein, Hagis; multiple authors]
\label{thm:omega}
If $n$ is an odd perfect number, then $\omega(n) \geq 9$,
where $\omega(n)$ denotes the number of distinct prime divisors of $n$.
\end{theorem}

The proof proceeds by showing that if $\omega(n) \leq 8$, the abundancy
constraint $\sigma(n)/n = 2$ cannot be satisfied.

\begin{theorem}[Chein 1979, Hagis 1980]
If additionally $3 \nmid n$, then $\omega(n) \geq 11$.
\end{theorem}

\begin{theorem}[Nielsen 2015]
$\omega(n) \geq 10$ if the largest component $p^a > n^{1/3}$.
\end{theorem}

\begin{remark}
Computational searches suggest that $\omega(n) \geq 10$ should hold in general,
and heuristic arguments point to $\omega(n)$ being quite large.
The exact minimum remains unknown.
\end{remark}

\subsection{Number of prime factors counted with multiplicity}

\begin{theorem}[Chein 1979]
If $n$ is an odd perfect number, then $\Omega(n) \geq 37$,
where $\Omega(n)$ counts prime factors with multiplicity.
\end{theorem}

This was later improved:

\begin{theorem}[Ochem--Rao 2014]
$\Omega(n) \geq 101$.
\end{theorem}

% ============================================================
\section{Bounds on Prime Factors}
% ============================================================

\subsection{The largest prime factor}

\begin{theorem}[Kanold, Hagis--McDaniel]
If $n$ is an odd perfect number, then its largest prime factor exceeds $10^8$.
\end{theorem}

More precisely:

\begin{theorem}[Goto--Ohno 2008]
The largest prime factor of an odd perfect number is greater than $10^8$.
\end{theorem}

\begin{theorem}[Luca--Pomerance 2010]
The second-largest prime factor of an odd perfect number exceeds $10^4$.
\end{theorem}

\subsection{The smallest prime factor}

\begin{theorem}
The smallest prime divisor of an odd perfect number is at least $3$
(trivially, since $n$ is odd).
\end{theorem}

\begin{theorem}[Grün 1952, extended]
If $3 \mid n$, then the smallest prime factor is $3$.
If $3 \nmid n$, the smallest prime factor is at least $5$.
\end{theorem}

\begin{theorem}[Iannucci 1999]
If $n$ is an odd perfect number with $3 \nmid n$ and $5 \nmid n$, then the
smallest prime factor of $n$ is at least $100003$.
\end{theorem}

% ============================================================
\section{Divisibility Constraints}
% ============================================================

\begin{theorem}[Various]
If $n$ is an odd perfect number, then:
\begin{enumerate}[label=(\roman*)]
  \item $n \equiv 1 \pmod{4}$ (Corollary~\ref{cor:form});
  \item $n \equiv 1 \pmod{12}$ or $n \equiv 9 \pmod{36}$ (Touchard 1953);
  \item $n$ is not divisible by $105 = 3 \cdot 5 \cdot 7$ (Steuerwald 1937);
  \item If $3 \mid n$, then $9 \mid n$ (various);
  \item $n$ cannot be a perfect square (direct from Euler's theorem,
        as $p^a$ with $a$ odd prevents this).
\end{enumerate}
\end{theorem}

\begin{theorem}[Touchard 1953]
\label{thm:touchard}
If $n$ is an odd perfect number, then either
$n \equiv 1 \pmod{12}$ or $n \equiv 9 \pmod{36}$.
\end{theorem}

\begin{proof}[Proof sketch]
This follows from Euler's theorem combined with a careful analysis of
$\sigma(n)/n = 2$ modulo $3$ and $4$ simultaneously.
\end{proof}

\begin{theorem}[Sevenster 2022]
If $3^2 \mid n$, the structure imposes additional congruence conditions
that severely restrict the possible factorisation types.
\end{theorem}

% ============================================================
\section{The Abundancy Index and Obstruction Theory}
% ============================================================

\begin{definition}
An \emph{abundancy obstruction} for a rational number $r$ is a proof that
no integer $n$ satisfies $\sigma(n)/n = r$.
\end{definition}

\begin{proposition}
For the equation $\sigma(n)/n = 2$ to hold with $n$ odd, the prime
factorisation of $n$ must be severely constrained: no sub-product of
prime-power components can create a ``local obstruction''.
\end{proposition}

\begin{theorem}[Broughan--Delbourgo--Zhou]
There is a dense set of rational numbers $r > 1$ for which no $n$ satisfies
$\sigma(n)/n = r$, but $2$ is \emph{not} among them based on current methods.
\end{theorem}

\begin{remark}
The non-existence of odd perfect numbers cannot be proved by local obstructions
alone: for every finite set of primes $S$, there exist odd numbers $n$ (not
perfect) whose $S$-local abundancy matches that of $2$.
\end{remark}

% ============================================================
\section{Related Concepts}
% ============================================================

\subsection{Quasiperfect numbers}

\begin{definition}
A positive integer $n$ is called \emph{quasiperfect} if $\sigma(n) = 2n + 1$.
\end{definition}

\begin{conjecture}
No quasiperfect numbers exist.
\end{conjecture}

It is known that any quasiperfect number must be an odd perfect square greater
than $10^{35}$ with at least seven distinct prime factors.

\subsection{Multiply perfect numbers}

\begin{definition}
A positive integer $n$ is \emph{$k$-perfect} (or \emph{multiply perfect of
index $k$}) if $\sigma(n) = kn$.
\end{definition}

For $k = 3$: $\sigma(n) = 3n$ (triperfect numbers; six are known:
$120, 672, 523776, 459818240, 1476304896, 51001180160$).
For $k \geq 4$ many examples are known but all are even.

\begin{conjecture}
For every $k \geq 3$, all $k$-perfect numbers are even.
\end{conjecture}

\subsection{Hyperperfect numbers}

\begin{definition}
A positive integer $n$ is \emph{$k$-hyperperfect} if
$n = 1 + k(s(n)) = 1 + k(\sigma(n) - n - 1)$, i.e.\ $\sigma(n) = (1+1/k)n + (1-1/k)$.
\end{definition}

\begin{remark}
For $k = 1$, hyperperfect reduces to perfect.
The existence of odd hyperperfect numbers for $k \geq 2$ is also open.
\end{remark}

% ============================================================
\section{Connections to Other Open Problems}
% ============================================================

\subsection{Goldbach--Lehmer connections}

Odd perfect numbers are related to Lehmer's totient problem: if $\phi(n) \mid (n-1)$
for composite $n$ (a Lehmer number), structural constraints resemble those for
odd perfect numbers.

\subsection{Abundant and primitive abundant numbers}

\begin{definition}
A positive integer $n$ is \emph{primitive abundant} if it is abundant but
no proper divisor of $n$ is abundant.
\end{definition}

Every even abundant number has a primitive abundant divisor.
The distribution of primitive abundant numbers intersects with constraints
on near-perfect numbers.

\subsection{Unitary perfect numbers}

\begin{definition}
$n$ is \emph{unitary perfect} if $\sigma^*(n) = 2n$, where
$\sigma^*(n) = \sum_{d \mid n,\, \gcd(d,n/d)=1} d$.
\end{definition}

Only five unitary perfect numbers are known; all are even.

% ============================================================
\section{Computational Evidence}
% ============================================================

\begin{theorem}[Ochem--Rao 2012, computational]
There is no odd perfect number less than $10^{1500}$.
\end{theorem}

Modern computational searches use:
\begin{enumerate}[label=(\arabic*)]
  \item \textbf{Factor chain methods}: enforcing $\sigma(p^a) \mid 2n$ for each
        prime power component.
  \item \textbf{Exhaustive case analysis}: by number of prime factors, size of
        Euler prime, etc.
  \item \textbf{Algebraic constraints}: combining Touchard's theorem,
        congruence conditions, and the Euler factorisation.
\end{enumerate}

\begin{remark}
The computational task of checking all odd numbers up to even $10^{100}$ for
perfection is entirely infeasible directly; the bounds are achieved by structural
arguments, not exhaustive search of all numbers.
\end{remark}

% ============================================================
\section{Heuristic Arguments}
% ============================================================

\subsection{Probability heuristic}

A standard heuristic suggests that the number of perfect numbers up to $x$
should grow very slowly.
For odd numbers, one argues as follows:
the probability that a ``random'' odd number $n$ satisfies $\sigma(n)/n = 2$
is related to the density of the set $\{r \in \Q : r = \sigma(n)/n\}$
near $2$.

\begin{remark}[Heuristic]
Since $\sigma(n)/n$ varies continuously as $n$ ranges over structured families,
and since the constraint is extremely rigid (exactly $2$, not approximately),
heuristic probabilistic arguments suggest the expected number of odd perfect
numbers is $0$, consistent with Conjecture~\ref{conj:nopn}.
\end{remark}

\subsection{The Euler prime constraint}

Because the Euler prime $p \equiv 1 \pmod{4}$ must appear to an odd power,
and because $\sigma(p^a)/p^a = (p^{a+1}-1)/((p-1)p^a)$ must together
with the remaining factors produce exactly $2$, the constraint becomes
increasingly rigid as $\omega(n)$ grows.

% ============================================================
\section{Known Necessary Conditions: A Summary}
% ============================================================

We collect all established necessary conditions:

\begin{theorem}[Necessary Conditions for Odd Perfect Numbers]
\label{thm:necessary}
If an odd perfect number $n$ exists, then:
\begin{enumerate}[label=(\roman*)]
  \item $n = p^a m^2$ with $p$ prime, $p \equiv a \equiv 1 \pmod{4}$, $\gcd(p,m)=1$;
  \item $n > 10^{1500}$;
  \item $\omega(n) \geq 9$;
  \item $\Omega(n) \geq 101$;
  \item The largest prime factor of $n$ exceeds $10^8$;
  \item The second-largest prime factor exceeds $10^4$;
  \item $n \equiv 1 \pmod{4}$;
  \item Either $n \equiv 1 \pmod{12}$ or $n \equiv 9 \pmod{36}$;
  \item $105 \nmid n$;
  \item $n$ is not a perfect square.
\end{enumerate}
\end{theorem}

\begin{remark}
Theorem~\ref{thm:necessary} represents the accumulated work of over two
centuries of number theory.
Each condition significantly narrows the search space, yet together they
still do not yield a contradiction.
\end{remark}

% ============================================================
\section{Conclusion and Open Questions}
% ============================================================

The problem of odd perfect numbers stands as one of the most ancient open
problems in all of mathematics, with no solution in sight despite the
combined force of analytic, algebraic, and computational methods.

The central conjecture remains:

\begin{conjecture}[Restated]
There are no odd perfect numbers.
\end{conjecture}

\noindent\textbf{Open questions of current interest include:}
\begin{enumerate}[label=(\arabic*)]
  \item Can $\omega(n) \geq 10$ be proved?
  \item Can the lower bound $n > 10^{1500}$ be improved substantially?
  \item Is there a prime $p$ that \emph{cannot} be the Euler prime of an odd
        perfect number?
  \item Can quasiperfect numbers be ruled out?
  \item Is there an algebraic number field argument that would force a
        contradiction in the structure $n = p^a m^2$?
  \item Could methods from analytic number theory (e.g.\ sieve methods) provide
        a non-computational approach?
\end{enumerate}

The problem connects deeply to the structure of the multiplicative function
$\sigma$ and to questions about the distribution of values of $\sigma(n)/n$
near $2$, touching on some of the most fundamental unsolved issues in
elementary number theory.

% ============================================================
\begin{thebibliography}{99}

\bibitem{Euler1747}
L.\ Euler,
\textit{De numeris amicabilibus},
Nova Acta Eruditorum (1747);
also in: Opera Omnia, Ser.\ I, Vol.\ 2, pp.\ 86--162.

\bibitem{Touchard1953}
J.\ Touchard,
\textit{On prime numbers and perfect numbers},
Scripta Mathematica \textbf{19} (1953), 35--39.

\bibitem{Kanold1941}
H.-J.\ Kanold,
\textit{Untersuchungen über ungerade vollkommene Zahlen},
J.\ Reine Angew.\ Math.\ \textbf{183} (1941), 98--109.

\bibitem{Chein1979}
E.\ Z.\ Chein,
\textit{An odd perfect number has at least 8 prime factors},
Ph.D.\ Thesis, Pennsylvania State University, 1979.

\bibitem{Hagis1980}
P.\ Hagis Jr.,
\textit{Outline of a proof that every odd perfect number has at least eight
prime factors},
Math.\ Comp.\ \textbf{35} (1980), 1027--1032.

\bibitem{BrentCohenRiele1991}
R.\ P.\ Brent, G.\ L.\ Cohen, H.\ J.\ J.\ te Riele,
\textit{Improved techniques for lower bounds for odd perfect numbers},
Math.\ Comp.\ \textbf{57} (1991), 857--868.

\bibitem{Iannucci1999}
D.\ E.\ Iannucci,
\textit{The third largest prime divisor of an odd perfect number exceeds
one hundred},
Math.\ Comp.\ \textbf{69} (2000), 867--879.

\bibitem{GotoOhno2008}
T.\ Goto and Y.\ Ohno,
\textit{Odd perfect numbers have a prime factor exceeding $10^8$},
Math.\ Comp.\ \textbf{77} (2008), 1859--1868.

\bibitem{LucaPomerance2010}
F.\ Luca and C.\ Pomerance,
\textit{On the radical of a perfect number},
New York J.\ Math.\ \textbf{16} (2010), 23--30.

\bibitem{OchemRao2012}
P.\ Ochem and M.\ Rao,
\textit{Odd perfect numbers are greater than $10^{1500}$},
Math.\ Comp.\ \textbf{81} (2012), 1869--1877.

\bibitem{OchemRao2014}
P.\ Ochem and M.\ Rao,
\textit{On the number of prime factors of an odd perfect number},
Math.\ Comp.\ \textbf{83} (2014), 2435--2439.

\bibitem{Nielsen2015}
P.\ P.\ Nielsen,
\textit{Odd perfect numbers, Diophantine equations, and upper bounds},
Math.\ Comp.\ \textbf{84} (2015), 2549--2567.

\bibitem{BroughanDelbourgoZhou}
K.\ A.\ Broughan, D.\ Delbourgo, Q.\ Zhou,
\textit{Improving the Chen and Chen result for odd perfect numbers},
Integers \textbf{13} (2013), A39.

\bibitem{Steuerwald1937}
R.\ Steuerwald,
\textit{Verschärfung eines Satzes über die Seite vollkommener Zahlen},
S.-B.\ math.-nat.\ Kl.\ Bayer.\ Akad.\ Wiss.\ 1937, 69--72.

\end{thebibliography}

\end{document}
